\section{Serializacja danych}
Aby przesłać dane za pomocą sygnału cyfrowego,
niezbędne jest przekształcenie ich na zapis bitowy.


\subsection*{Dane pomiarowe}
Do interpretacji danych pomiarowych niezbędne są 3~informacje:
\begin{itemize}
    \item \textbf{typ danych} (SENSOR TYPE) --- z~jakiego typu urządzenia
    pomiarowego pochodzą dane (np. termometr, barometr itp.)
    \item \textbf{numer urządzenia pomiarowego} (ID) --- czujnik 
    klimatyczny może posiadać wiele modułów pomiarowych
    tego samego typu.
    \item \textbf{rozdzielczość danych} (SCALE) --- ilość liczb po
    po przecinku zawartych w~pomiarze.
\end{itemize}
Ze względu na chęć zachowania jak najmniejszego rozmiaru,
informacje te znajdą się w~jednym bajcie poprzedzającym 
samą odczytaną wartość. Bajt ten posiada następującą
strukturę:

\bigskip

\begin{bytefield}[bitwidth=4.4em,endianness=big]{8}
    \bitheader{0,3,4,5,6,7} \\
    \bitbox{2}{ID}
    \bitbox{2}{SCALE}
    \bitbox{4}{SENSOR TYPE}
\end{bytefield}

Wartość pomiaru zapisywana jest jako 32~bitowa liczba bez 
znaku. W~przypadku danych, w~których występują liczby
ujemne (np. temperatura), konieczna jest konwersja na liczbę
bez znaku podczas jej serializacji i~następnie konwersja
na liczbę ze znakiem przy deserializacji.
Typ ten został wybrany z~powodu możliwości zapisania
w~nim danych z~każdego obsługiwanego w~projekcie modułu.

Całość wartości pomiarowej wraz z~danymi ją opisującymi 
ma rozmiar 5 bajtów i~strukturę:

\bigskip

\begin{bytefield}[endianness=little,bitwidth=0.8em]{40}
    \bitheader{0,7,8,39} \\
    \bitbox{8}{LABEL \\ \tiny 1 byte} 
    \bitbox{32}{DATA \\ \tiny 4 bytes}
\end{bytefield}

\subsection*{Serializacja w kodzie}
Za serializację danych pomiarowych odpowiedzialna 
jest struktura \texttt{TypedLabelReadout}, przedstawiona na 
listingu~\ref{lst:typed_lebel_readout}.
Posiada ona ona czterobajtowe pole danych typu \texttt{u32} 
oraz trzy parametry generyczne \texttt{ID}, \texttt{SCALE}
i~\texttt{TYPE}, odpowiedzialne odpowiednio za: numer urządzenia
pomiarowego, rozdzielczość danych i~typ urządzenia pomiarowego.

\begin{lstlisting}[
    language=Rust, style=colouredRust, 
    caption={Struktura \texttt{Transmitter}},
    label={lst:typed_lebel_readout}
]
pub struct TypedLabelReadout<ID, SCALE, TYPE>{
    data: u32,
    _label: PhantomData<(ID, SCALE, TYPE)>,
}
\end{lstlisting}